% notation-table.tex — 附录：符号总表

\chapter{符号总表}\label{app:notation}

本附录汇总全书使用的所有符号，便于快速查阅。详细定义见 \cref{ch:notation}。

\section{图像与分量}

\begin{tabular}{@{}ll@{}}
\toprule
符号 & 含义 \\
\midrule
$W$ & 图像宽度（像素） \\
$H$ & 图像高度（像素） \\
$C$ & 分量总数 \\
$c$ & 分量索引，$c = 0, 1, \ldots, C-1$ \\
$(x, y)$ & 像素坐标 \\
$d$ & 样本位深（bit depth） \\
$s_x[c], s_y[c]$ & 分量 $c$ 的水平/垂直子采样因子 \\
\bottomrule
\end{tabular}

\section{内部精度与 NLT 参数}

\begin{tabular}{@{}ll@{}}
\toprule
符号 & 含义 \\
\midrule
$B_w$ & 内部工作位宽 \\
$F_q$ & 小数位数（fractional bits） \\
$T_1, T_2$ & NLT Extended 阈值 \\
$E$ & NLT Extended 线性区斜率控制 \\
$\sigma, \alpha$ & NLT Quadratic 参数 \\
\bottomrule
\end{tabular}

\section{小波分解}

\begin{tabular}{@{}ll@{}}
\toprule
符号 & 含义 \\
\midrule
$NL_x$ & 水平分解层数 \\
$NL_y$ & 垂直分解层数 \\
$b$ & 子带索引 \\
$S_d$ & 小波变换抑制的分量数 \\
\bottomrule
\end{tabular}

\section{Precinct 与 Packet}

\begin{tabular}{@{}ll@{}}
\toprule
符号 & 含义 \\
\midrule
$C_w$ & 列宽（column width），以 8 为单位 \\
$H_{sl}$ & slice 高度（以 precinct 行为单位） \\
$(p_x, p_y)$ & precinct 列、行索引 \\
$N_g$ & GCLI 组大小（group size） \\
$S_s$ & 显著性标志步长 \\
\bottomrule
\end{tabular}

\section{量化与码率控制}

\begin{tabular}{@{}ll@{}}
\toprule
符号 & 含义 \\
\midrule
$GCLI_b$ & 子带 $b$ 的全局 CLI \\
$GTLI_b$ & 子带 $b$ 的全局截断位位置 \\
$Q$ & 量化参数（quantization） \\
$R$ & 细化参数（refinement） \\
$B_{total}$ & 总预算（bits） \\
$B_{cbr}$ & CBR 分配预算 \\
$B_{prec}$ & precinct 预算 \\
\bottomrule
\end{tabular}
